\begin{table}[h]
\centering
\caption{Common Categories of Knowledge Graph Recommender Systems \parencite{yu2024review}.}
\begin{tabular}{|p{0.35\textwidth}|p{0.55\textwidth}|}
\hline
\textbf{Typology} & \textbf{Clarification} \\
\hline
Recommender system based on graph neural network & Such systems model the nodes and edges of the knowledge graph through graph neural networks, learn the representation of users and items, and implement recommendations by predicting ratings or preferences. \\
\hline
Path Search Based Recommender System & This type of system mines items that may be of interest to users by searching for paths in the knowledge graph, e.g., searching for items that similar users like based on their historical behavior and preferences. \\
\hline
Recommender systems based on graph embedding & This type of system maps the nodes and edges of the knowledge graph to a low-dimensional vector space, capturing relationship and attribute information, and utilizes the low-dimensional vectors to make recommendations, e.g., by calculating the similarity between users and items to make recommendations. \\
\hline
Recommender systems based on heterogeneous data from multiple sources & This type of system integrates multi-domain data and uses knowledge graphs for data representation and processing to achieve cross-domain recommendations, such as combining social media behavior and e-commerce purchase records for more comprehensive recommendations. \\
\hline
Recommender systems based on multi-task learning & Such systems combine multiple related recommendation tasks and use knowledge graph information for collaborative learning, e.g., simultaneous item recommendation and user label prediction, to improve recommendation accuracy and personalization. \\
\hline
\end{tabular}
\end{table}